\documentclass[11pt,a4paper]{article}
\usepackage[margin=23mm]{geometry}
\usepackage{amsmath,amssymb,amsthm}
\setlength{\emergencystretch}{2em}
\newtheorem{lemma}{Lemma}
\newtheorem{theorem}[lemma]{Theorem}
\newcommand{\gen}[1]{\langle#1\rangle}
\newcommand{\Sym}{\operatorname{Sym}}
\title{The witness ending one step after the crossing}
\author{Human Compression project: L044 continuation}
\date{}
\begin{document}\maketitle
All marks are involutions satisfying string commutations.
The local actions below need no intersection property.

\begin{lemma}[One right mark]
Let $Q=\gen{\rho_m,\ldots,\rho_{j+1}}$ be transitive
on $W$, with every displayed mark active. Suppose the
full local $j$-omission has two orbits joined by exactly
one $j$-edge, the first omission $Q_{>m}$ is
intransitive, and
\[
                    h=j-m\ge0,\qquad |W|=h+3.
\]
Then $Q_{>m}$ has precisely two orbits, a singleton
and its transitive complement.
\end{lemma}
\begin{proof}
Put $L=Q_{<j}$ and $R=\gen{\rho_{j+1}}$. The
actual omission sides $A,B$ are transitive under
$LR$, with crossing endpoints $a,b$.

First suppose $h\ge3$. The sharper acquired
commuting-factor bound excludes attaining the
left length $h$ and right length one on the
same side: that side alone would require $h+3$
points. Thus, after naming the sides,
\[
 |A|=h+1,\quad |B|=2,\quad L|_A=\Sym(A),\quad
 R|_A=1,\quad R|_B=\Sym(B).
                                                        \tag{1}
\]
The same conclusion holds for $h=2$, unless both
maxima lie on a four-point side $A$, with the
other side a singleton. We exclude that exception
explicitly. The single involution $R$ cannot be
transitive on $A$, and the earlier commuting-factor
count would give six points if neither factor were
transitive. Hence $L$ is transitive on $A$. Write
$t_0=\rho_{j-1}|_A$, $t_1=\rho_m|_A$, and
$c=\rho_{j+1}|_A$. The involution $c$ is a nontrivial
centralizer of the transitive $L$, hence is fixed-point-free.
The mark $t_1$ fixes $a$, and therefore also $a^c$.
Being nonidentity, it is the transposition of the
other two points. Transitivity of $L$ forces
$t_0$ to exchange the two $c$-pairs. Commutation
with $c$ then makes $t_0$ a product of two
cross-pair transpositions. Consequently
$\gen{t_0,c}$ is transitive on the four points.
These marks survive the first omission, as does
the crossing to the singleton. The first omission
would be transitive, a contradiction. This proves
(1) for $h=2$ also.

Under (1), every left mark except $\rho_{j-1}$
fixes $b$, because it commutes with $\rho_j$ and
preserves the unique crossing. It thus fixes
both points of $B$. The nearest left mark may
also transpose $B$, but $R$ itself supplies that
transposition. On $A$, the exhausted anchored
path is $a_0=a,\ldots,a_h$, in reverse left
order. The only possible extra path transposition
is the acquired decoration $(a_{h-1}\ a_h)$ in
the mark three positions from the end. If present,
the remaining left marks after omitting $\rho_m$
would still connect all of $A$. Together with
$R$ and the crossing they would make the first
omission transitive. Hence that decoration is absent,
and its remaining left marks connect exactly
$A\setminus\{a_h\}$ and fix $a_h$.

For $h\ge3$, the far left marks give the pure
$\Sym(\{a_1,\ldots,a_h\})$, fixing $B$.
Commuting with this group makes $\rho_j$ fix
its entire support. For $h=2$, the only possible
additional action of $\rho_j$ on $A$ is
$(a_1\ a_2)$; it would connect the first omission
and is therefore absent. Thus $z=a_h$ is fixed
by the first omission, whose other points are
connected by the remaining left action, the
crossing and $R$. This gives the assertion for
all $h\ge2$.

If $h=1$, there are four points. An omission
side has transitive action of the commuting
involutions $L,R$, so its size is a power of two.
Thus the sides each have size two. The unique
crossing is the literal transposition of their
chosen endpoints. If $R$ transposed both sides,
$\gen{\rho_j,R}$ would be transitive on all four
points, contrary to the first-omission hypothesis.
It must transpose exactly one side. Its crossing
therefore gives a transitive three-point orbit
and fixes the unused point of the other side.
If $h=0$, there are three points, $L=1$, and
the active involution $R$ is a transposition.
It is itself the first omission, proving the
last case.
\end{proof}

\begin{lemma}[An anchor in the left endpoint orbits]
Let $Q=\gen{\rho_m,\ldots,\rho_y}$ be transitive
on $W$, every mark active, with a unique full
local $j$-crossing $\{a,b\}$ and $y\ge j+2$.
Put $h=j-m\ge0$, $L=Q_{<j}$, $R=Q_{>j}$ and
\[
                      C=a^L\sqcup b^L.
\]
Suppose $|C|=h+2$ and some $z\in C$ is fixed
by the first omission $Q_{>m}$. Then that omission
has precisely the orbits $\{z\}$ and $W\setminus\{z\}$.
\end{lemma}
\begin{proof}
For $h\ge1$, the anchored bound and equality
make one of these two $L$-orbits a set $A_0$
of size $h+1$, with $L|_{A_0}=\Sym(A_0)$,
and the other a singleton $b$, after renaming.
Since $\rho_j$ belongs to the first omission
and exchanges $a,b$, the fixed point $z$ lies
in $A_0\setminus\{a\}$. The whole right factor
fixes $z$ and centralizes $L$, so fixes every
point of $A_0=z^L$. Thus the complete actual
omission side through $a$ is $A=A_0$.
The left factor fixes $b$ and centralizes $R$,
so the other actual side is $B=b^R$ and
$L$ fixes all of $B$.

For $h\ge2$, the exhausted anchored path
$a_0=a,\ldots,a_h$ shows that the surviving
left marks connect $a_0$ through $a_{h-1}$.
Thus their fixed point $z\ne a_0$ can only be
$a_h$; a possible extra edge joining $a_h$
is excluded by the assumed fixation of $z$.
For $h=1$, the set $A\setminus\{z\}$ is just
$\{a\}$. In every case the surviving left
marks are transitive on $A\setminus\{z\}$.
The right factor is transitive on $B$, and
$\rho_j$ joins these sets while fixing $z$.
This proves the conclusion.

If $h=0$, then $C=\{a,b\}$ and the first
omission is $R$. Its fixed anchor $z$ is one
of $a,b$. The actual omission side through
$z$ is therefore a singleton; $R$ is transitive
on the other side. This proves this case too.
\end{proof}

Retain the complete actual odd-degree head and
paired-tail hypotheses, and the strong simultaneous
choice of crossing colour $j\ge5$. Write
$H=G_{\ge5}$, $K=G_{>5}$, $\sigma=\rho_4$.
Every $K$-intransitive witness orbit has first
active label $5$, a unique full local crossing,
and pointwise fixation by $\delta$ in the actual
head interface $\rho_3=(6\ b)\delta$.

\begin{theorem}
Let $O_1,O_2$ be two distinct $K$-intransitive
witness orbits, with last active labels $x=j+1<y$.
Then $x<r-1$ and
\[
                   |\operatorname{Fix}(G_{>x})|\ge2x+1.
\]
\end{theorem}
\begin{proof}
The acquired simultaneous-witness count gives
$x<r-1$ and at least $2x$ fixed points. If the
asserted stronger bound fails, the fixed set
$F=\operatorname{Fix}(G_{>x})$ is precisely
\[
 F=D\cup\{b\}\ \sqcup\ O_1\ \sqcup\ C,
 \qquad |O_1|=x-3,\quad |C|=x-4,                 \tag{2}
\]
where $C$ is the union of the two left endpoint
orbits in $O_2$. In particular $|C|=h+2$ for
$h=j-5$, and the first lemma applies on $O_1$.
Its $K$-orbits are a singleton $\{z_1\}$ and
a transitive complement $T_1$, with $|T_1|=|C|$.

The actual anchor $z_0=b^\sigma$ is outside
the seven head points and is fixed by $K$.
It lies in $F$, so (2) places it either at
$z_1$ or in $C$. The involution $\sigma$
centralizes $K$, hence permutes its actual
orbits. It also preserves $F$, since $x\ge6$.

Whenever the $K$-orbits of a witness $O_i$
are $\{z_0\},T_i$, the involution $\sigma$
cannot preserve $T_i$. Otherwise the proper
set $D\cup\{b\}\cup O_i$ would be invariant
under every generator: the first three marks
preserve $D$; $\rho_3$ exchanges $6,b$ and
fixes $O_i$ through $\delta$; $\sigma$
exchanges $b,z_0$, fixes $D$ and preserves
$T_i$; and $H$ preserves $O_i$ and fixes the
seven head points. The other witness is outside
this set, contrary to transitivity of $G$.

Suppose first $z_0=z_1$. The distinct $K$-orbit
$T_1^\sigma$ lies in $F$. It cannot meet the
fixed head points or $z_1$, and cannot equal
$T_1$. Equation (2) therefore puts it inside
$C$. Their equal cardinalities give
$T_1^\sigma=C$, so $C$ is $K$-invariant.
It is also invariant under $\rho_5$: for
$j>5$ this mark belongs to $L$, whose two
orbits define $C$; for $j=5$, $C$ consists
of the two crossing endpoints, interchanged
by $\rho_5$. Thus $C$ is $H$-invariant.
It must equal the whole actual orbit $O_2$.
But $C\subset F$ and $\rho_y$ is active on
$O_2$, with $y>x$, a contradiction.

Suppose instead $z_0\in C$. The second lemma
makes the $K$-orbits on $O_2$ exactly
$\{z_0\},T_2$. By the preceding invariant-set
argument $\sigma$ moves $T_2$ to a distinct
$K$-orbit. Take any point of
$C\setminus\{z_0\}$; this set is nonempty,
since $|C|=x-4\ge2$. Its image is in $F$
and in the nontrivial orbit $T_2^\sigma$.
It cannot be one of the $K$-fixed points,
and cannot lie in $O_2$, whose only nontrivial
$K$-orbit is $T_2$. By (2) it lies in $T_1$,
so $T_2^\sigma=T_1$. This is impossible:
$\rho_y\in K$ is nontrivial on $T_2$ but
fixes $T_1$ pointwise, while $\sigma$
centralizes it. Both cases contradict the
assumption $|F|=2x$, proving the bound.
\end{proof}

\paragraph{Attribution and boundary.}
The actual head interface, simultaneous witness
selection, anchored paths and orbit counts are
acquired. The two local equality arguments and
the action of the head anchor on the entire
fixed set are supplied here. All short cases
are included. The case $x=j$, a second witness
with transitive $K$-action and full outside
attachments remain. No whole-original, question,
formal or numerical credit is claimed.
\end{document}
